\note{sa page 10}

Or si $i \leqslant n$, on aura

\[ (6.12) \qquad \begin{cases} \Lambda\, \pi_{i,0} = 0 , \\[2pt] \Lambda\, \pi_{i,\alpha} = \struck{L}\ w_{i-2\alpha}\, L^{n-i+\alpha}\, \pi_{i,\alpha} \quad \text{si } \alpha \geqslant 1 , \end{cases} \]

\struck{et si $j > n$, \ill{} (6.13) $\Lambda \pi_{j,\alpha} = L \ill{}$}

\[ (6.13) \qquad \Lambda\, \pi_{2n-i,\alpha} = L^{n-i+\alpha-1}\, w_{i-2\alpha}\, \struck{\ill{}}\ L^{\alpha}\, \pi_{2n-i,\alpha} \qquad (\text{si } i \leqslant n-1) \]

Les conditions envisagées indiquent \ill{} (iv), i.e.\ prouvons que l'on a
alors $\pi_{i} = \pi'_{i}$. (iv) indiquent (ii). Il suffit de \ill{}
\emph{Démonstration à fournir.}

\emph{Remarque.} Voici une \ill{} variante de \ill{} unique \ill{} en termes
de $L$, et \ill{} une caractérisation \ill{} en termes de l'anneau gradué
$\mathcal{E}$) :

\marginal{$\Lambda \in \mathcal{E}^{(2)}$, $\pi_{i,\alpha} \in \mathcal{E}^{(0)}$}

\note{En marge il écrit $\mathcal{E}^{(2)}$ pour $\Lambda$, alors que partout
ailleurs $\Lambda$ est de degré $-2$ ; le signe manque sur le feuillet.}

\begin{enumerate}
  \item $\bigl( \Lambda^{n-i} L^{n-i} - 1 \bigr)\, \pi_{i} = 0$ ;
  \item $\bigl( L^{n-i} \Lambda^{n-i} - 1 \bigr)\, \pi_{2n-i} = 0$ ;
  \item $\pi_{i,\alpha} = L^{\alpha}\, \Lambda^{n-i+2\alpha}\, L^{n-i+\alpha} \Bigl( 1 - \sum_{\alpha' > \alpha} \pi_{i,\alpha'} \Bigr)$ ;
  \item $\pi_{2n-i,\alpha} = L^{n-i}\, \pi_{i,\alpha}\, \Lambda^{n-i}$ .
\end{enumerate}

